\paragraph{The coarse Weil estimate still does not supply A47.}
This is a limitation of a proposed estimate, not a lower bound
for the actual exponent. Work in genus two and characteristic
two. First suppose that the cyclic deck group of a pointed
cover of degree $n>1$ is constant over $\mathbf F_q$, and put
\[
 B=(\sqrt q+1)^2,\qquad f=\operatorname{ord}_n(2),
                         \qquad r=\varphi(n)/f\ge1.
 \tag{BPW.1}
\]
The pointed abelian-cover classification, equivalently the
Abel-map fundamental-group statement and the finite-field
Lang quotient, makes its cyclic group a quotient of
$J_C(\mathbf F_q)$ \cite[Proposition~9.1]{MilneJV}.
The Weil bound therefore gives
\[
                  n\mid\#J_C(\mathbf F_q),\qquad n\le B^2
 \tag{BPW.2}
\]
\cite[Theorem~19.1]{MilneAV}.

The primitive factor has dimension $\varphi(n)$: the
degree-$d$ quotient has genus $d+1$, and the cyclic conductor
decomposition of its Prym dimensions is
$d-1=\sum_{a\mid d,\ a>1}\varphi(a)$.
Since the deck maps are rational, $P_n(\mathbf F_q)[2^\infty]$
is a module over
$\mathcal O_n=\mathbf Z_2[Z]/\Phi_n(Z)$.
This is a product of unramified discrete valuation rings of
common residue degree $f$. If a component contains an element
of order $M_n=2^a$, its cyclic module has annihilator exactly
$(2^a)$ and cardinality $2^{fa}=M_n^f$.
Taking a component of maximal exponent and applying the Weil
bound to the actual primitive abelian subvariety yields
\[
 M_n^f\le\#P_n(\mathbf F_q)\le B^{\varphi(n)},
                         \qquad M_n\le B^r.
 \tag{BPW.3}
\]
The unrounded substitution of (BPW.3) in (BP.2) would require
$n>(2^{\omega(n)}B^r-1)^2$. Its right side exceeds $B^2$,
since $2^{\omega(n)}\ge2$, $r\ge1$, and $B>1$.
It is incompatible with (BPW.2).

Rounding down to the largest power of two allowed by (BPW.3)
does not rescue this estimate. For $q\ge8$, that rounded
upper bound is at least $q$, because $q$ is itself a power
of two and $B^r\ge B>q$. But
\[
           2q-1-B=q-2\sqrt q-2>0\qquad(q\ge8).
 \tag{BPW.4}
\]
The inequality holds at $q=8$ since $6>4\sqrt2$, and the
expression increases for $q\ge8$. For $q=2$, the largest
power of two at most $B=3+2\sqrt2$ is four, and
$2\cdot4-1=7>B$. For $q=4$, it is eight since $B=9$,
and $2\cdot8-1=15>B$.
For larger $r$ the rounded upper bound can only increase.
Thus even the rounded substitution in (BP.2) always asks for
$n>B^2$, contradicting the available bound.

For a nonconstant geometric deck group, first pass to its own
minimal deck field. The actual primitive exponent is unchanged
by (SD.10); the pointed cover then satisfies the constant-deck
hypotheses just used. No full cyclotomic action on the rational
group before that passage is assumed. The same estimate fails
over the deck field as well.

This leaves the exact finite-field primitive exponent as a
missing input. Geometric rank information or local character
values can make it much smaller than (BPW.3); the argument
does not exclude such covers. Hooley's candidate-order density
does not supply that additional information. This supplements
A34's earlier coarse-bound obstruction and is not a new
universal nonexistence or main-route supply result.
